%% =====================================================================
\section{Introduction}
\label{sec:intro}
%% =====================================================================

The availability of a web service is as much a security property as its
confidentiality or integrity. Denial-of-service (DoS) attacks remain among the most
common and most damaging threats to online systems, and the research community has
produced extensive taxonomies and defenses for network- and transport-layer flooding
\cite{mirkovic2004,zargar2013}. Large botnets such as Mirai have shown that raw
volumetric capacity is cheap for an adversary to obtain \cite{antonakakis2017}. A
harder class of attacks, however, operates at the application layer: the adversary
issues syntactically and semantically valid requests that are individually
indistinguishable from those of legitimate users, so that volume- or signature-based
filters are ineffective or must be tuned so aggressively that they harm real customers
\cite{praseed2019,xie2009,ranjan2009}. This paper studies a particularly subtle member
of that class in the setting in which most modern software is now deployed: the
containerized, multi-service web application.

\subsection{Motivation: the Cache as an Attack Surface}
Virtually every read-heavy web application interposes an in-memory cache such as
Redis or memcached between its application tier and its database, so that repeated
requests for the same object are answered from memory rather than recomputed
\cite{fitzpatrick2004,nishtala2013}. Each cached entry receives a time-to-live (TTL),
after which it expires and must be regenerated from the database. The cache is usually
discussed as a \emph{performance} component, but it also changes the \emph{capacity
model} of the system: once a cache is in place the database is typically provisioned
only for the residual miss traffic, not for the full request rate. Production
workloads are highly skewed towards a small set of popular (``hot'') keys
\cite{breslau1999,atikoglu2012}, so the protection the cache provides is concentrated
on exactly those keys.

If a hot key expires while many requests for it are in flight, all of them miss the
cache at the same instant and are forwarded to the database simultaneously. A single
expiry therefore produces not one query but $N$ redundant queries, where $N$ is the
instantaneous concurrency on that key. This is a \emph{cache stampede}, also called a
\emph{dog-pile} or \emph{thundering herd}, and it is documented as a first-order hazard
in large-scale systems \cite{vattani2015,nishtala2013}. The consequences can be
severe even without an adversary: Facebook's widely discussed outage of September 2010
was attributed to an automated system that, on seeing an invalid cached configuration
value, caused every client to query the backing database cluster at once, overwhelming
it in a self-reinforcing loop \cite{johnson2010}. Practitioner guidance on cascading
failure likewise identifies the loss of a cache layer as one of the classic triggers of
system-wide overload \cite{beyer2016}.

What is less often treated is that the same mechanism is a \emph{security}
vulnerability. An adversary who understands a target's caching behavior can
deliberately align bursts of requests for hot objects with those objects' expiry,
repeatedly, so that the database is driven into overload by traffic that never leaves
the normal request distribution of the application. The database---not the web
server---is the victim, and the web application is merely the conduit. We call this
an \emph{Indirect Distributed Denial of Service} (IDDoS): \emph{indirect} because the
attacker never addresses the victim component, and \emph{distributed} in its effect
because the damage is produced by the fan-out of concurrency inside the system rather
than by the attacker's own request volume. Because the attack rides on correct
application logic rather than on a software defect, it cannot be ``patched''; it must
be \emph{detected} and \emph{mitigated} as an operational property of the system, in
the same way that algorithmic-complexity attacks exploit legitimate but expensive code
paths \cite{crosby2003} and low-rate attacks exploit legitimate protocol timers
\cite{kuzmanovic2003}.

\subsection{Problem Statement and the Complex Engineering Challenge}
Defending this vulnerability in a containerized environment is a genuine complex
engineering problem: its difficulty arises from the \emph{interaction} of
well-understood components rather than from any single exotic one. Three tensions
must be resolved simultaneously.

\begin{enumerate}[label=\textbf{T\arabic*.},leftmargin=*]
\item \textbf{Detection without a volumetric signature.} The attack's request rate at
the web tier can resemble a legitimate spike or flash crowd \cite{jung2002,kandula2005}.
A detector keyed on request volume alone will either miss the attack or raise
constant false alarms on benign bursts that the cache absorbs harmlessly.
\item \textbf{Mitigation without self-inflicted denial of service.} A defense that
simply rejects traffic to protect the database can itself become the denial of
service. The mitigation must shed the \emph{redundant} load while continuing to serve
genuine users quickly.
\item \textbf{Faithful, safe, repeatable experimentation.} Studying an availability
attack requires driving a realistic system into failure repeatedly and measuring
per-component behavior, which demands an isolated, resettable, production-like
environment rather than a single-process simulation---and one that cannot harm any
third party.
\end{enumerate}

These requirements are coupled. An overly aggressive guardrail degrades legitimate
users; a cache-only fix that ignores the connection path leaves users slow even when
the database is healthy (Section~\ref{sec:ablation}); and a detector that is evaluated
only on the traffic it was tuned on can look perfect while measuring the wrong thing
(Section~\ref{sec:det-results}). Resolving the problem therefore requires an integrated
architectural, algorithmic, and experimental design, not merely a faster database or a
larger cache.

\subsection{Research Questions}
The work is organized around four research questions:
\begin{description}[leftmargin=1.2em,style=nextline]
\item[RQ1 (Reproduction).] Can a cache-stampede IDDoS be reproduced deterministically
on an ordinary containerized stack using only well-formed requests, and what are its
quantitative effects on database load, latency, and availability?
\item[RQ2 (Mechanism).] Which mechanism converts the stampede into user-visible failure,
and can its magnitude be predicted from simple system parameters (key count, TTL,
request rate, connection-pool size)?
\item[RQ3 (Detection).] How do an explainable multi-signal rule and an unsupervised
IsolationForest compare in detection latency, coverage, false alarms, and
interpretability---and how do they compare with naive single-signal detectors?
\item[RQ4 (Mitigation).] Which combination of cache locking, stale-while-revalidate,
and rate limiting restores availability, how quickly, and what does each layer
contribute?
\end{description}

\subsection{Proposed Solution}
We present \textbf{CacheCollapse}, a framework that addresses all three tensions
together:
(i)~a \emph{reproducible containerized testbed} that reconstructs the
Nginx~$\rightarrow$~Flask~$\rightarrow$~Redis~$\rightarrow$~PostgreSQL request path on
an isolated Docker Compose network \cite{merkel2014,dockercompose} and generates both
normal traffic and a controllable stampede attack;
(ii)~a \emph{dual detection subsystem} pairing an explainable multi-signal rule with
an unsupervised IsolationForest \cite{liu2008,liu2012} trained only on baseline
behavior; and
(iii)~a \emph{layered mitigation} combining single-flight cache locking with
stale-while-revalidate \cite{nishtala2013,rfc5861}, an application-layer per-client
rate limiter, and an edge rate limiter at the reverse proxy \cite{nginxlimitreq}, each
of which can be toggled live from a browser dashboard while the attack is running.

\subsection{Contributions}
\begin{itemize}[leftmargin=*]
\item A \textbf{containerized IDDoS testbed} that reproduces the cache-stampede attack
on nine interacting industry-standard services, with a Redis-based control plane that
switches traffic, attack, and each defense on and off live
(Section~\ref{sec:design}).
\item An \textbf{analytical model} of baseline miss traffic, stampede amplification,
and defended load that predicts the measured baseline database rate to within 9\%
and explains why, under overload, completed database work collapses far below
nominal pool capacity (Sections~\ref{sec:model} and~\ref{sec:mechanics}).
\item A \textbf{multi-signal rule-based detector} and an \textbf{unsupervised
IsolationForest detector}, and a detailed comparison on 12 attack episodes and
${\approx}1{,}100$ labelled intervals covering detection latency, interval-level
precision and recall, behavior on idle and defended traffic, threshold sensitivity,
ROC analysis, and training-window sensitivity (Section~\ref{sec:det-results}).
\item A \textbf{layered mitigation}, an \textbf{ablation}, and a \textbf{longitudinal
analysis} of 14 experiment sessions showing how each engineering fix changed
defended-mode latency and request admission, including the deployment-critical
finding that database-level and proxy-level protection are distinct and both
required (Sections~\ref{sec:mitigation-results}--\ref{sec:ablation}).
\item A \textbf{fully reproducible, instrumented deployment} (single
\texttt{docker compose up}, Prometheus/Grafana, live dashboard, demonstration
storefront) validated by a 16-case automated end-to-end test suite, plus an explicit
account of \textbf{measurement pitfalls} in the instrumentation that other
practitioners are likely to repeat (Section~\ref{sec:validity}).
\end{itemize}

\subsection{Paper Organization}
Section~\ref{sec:background} reviews the technical background on caching, overload,
rate limiting, and anomaly detection. Section~\ref{sec:related} surveys related work
and identifies the research gap. Section~\ref{sec:threat} defines the threat model and
develops an analytical model of the attack. Section~\ref{sec:design} describes the
design and implementation of CacheCollapse. Section~\ref{sec:method} sets out the
experimental methodology and dataset. Section~\ref{sec:results} presents the results
for each research question. Section~\ref{sec:comparative} compares CacheCollapse with
prior work and alternative mitigations. Section~\ref{sec:discussion} discusses
implications, engineering lessons, threats to validity, and limitations, and
Section~\ref{sec:conclusion} concludes. Appendices give a reproduction guide, key
configuration listings, and a mapping of the project to the attributes of a complex
engineering problem.
